% 17-interop.tex — Clause 17: C interoperability
\chapter{C interoperability}
\label{sec:17-interop}
\index{C interoperability}

\section{Overview}
\label{sec:interop.overview}

\pnum
Lain compiles to C99 and provides direct, zero-overhead interoperability
with C libraries through three mechanisms:
\begin{enumerate}
  \item \lain{c\_include} --- injects C \texttt{\#include} directives;
  \item \lain{extern} declarations --- declares C functions for use in Lain;
  \item type mapping --- Lain types map to their C99 counterparts.
\end{enumerate}

\section{The \texttt{c\_include} declaration}
\label{sec:interop.cinclude}
\index{c\_include}

\pnum
\lain{c\_include "header.h"} causes the compiler to emit a
\texttt{\#include} directive in the generated C code.  The string content
is emitted verbatim.

\begin{laincode}
c_include "<stdio.h>"
c_include "<stdlib.h>"
\end{laincode}

\begin{constraint}
  \lain{c\_include} shall appear only at module top level.
\end{constraint}

\section{Extern function declarations}
\label{sec:interop.extern}
\index{extern function}

\pnum
\lain{extern func name(params) return} declares a C function.  Because an
extern has no body, its effect row cannot be inferred; the row is therefore
\emph{believed}, and its default is the set of \emph{all} effects
(\lain{io}, \lain{diverge}, \lain{raises}, \lain{alloc}).  An
\lain{effects} clause on an extern \emph{narrows} that default.

\pnum
This is the inverse of a declaration with a body, whose default row is
$\emptyset$ and whose clause \emph{widens} it~(\S\,\ref{sec:decl.func}).
The direction differs for one stated reason: inference is available in the
one case and impossible in the other.  A caller of an extern acknowledges
whatever the extern's row names, by the ordinary rule.

\begin{constraint}
  An \lain{extern} with no \lain{effects} clause has the row
  \{\lain{io}, \lain{diverge}, \lain{raises}, \lain{alloc}\}.  A caller
  whose own bound does not cover it is \illformed{}~(\diagcode{E011}).
\end{constraint}

\begin{note}
  \textbf{Trust boundary.}
  An extern's \lain{effects} clause is a \textit{programmer assertion} that
  the referenced C function does no more than the row names.  The compiler
  cannot verify it and believes it, so a wrong narrowing---declaring
  \lain{libc\_rand} with an empty row, say---may invalidate the guarantees
  that depend on purity (notably P4 determinism for any function that
  transitively calls it).  What has changed is the cost of \emph{silence}:
  omitting the clause once meant "trusted pure" (the rule once labelled I-016), so
  forgetting it on a C function that performs I/O was a silent fail-open.
  Omitting it now claims everything, so forgetting costs precision instead.
  The empty row, \lain{effects} with no names, is how purity is asserted
  deliberately.
\end{note}

\pnum
Ownership annotations on extern parameters work identically to those on
native Lain functions; the borrow checker tracks ownership of resources
obtained via extern calls.

\section{Type mapping}
\label{sec:interop.types}
\index{type mapping}

\pnum
The following table defines the normative mapping between Lain types and
their C99 equivalents.  The mapping is the ABI used by the translation
output.

\begin{longtable}{@{}lll@{}}
\toprule
\textbf{Lain type} & \textbf{C99 type} & \textbf{Notes} \\
\midrule
\endfirsthead
\toprule
\textbf{Lain type} & \textbf{C99 type} & \textbf{Notes} \\
\midrule
\endhead
\bottomrule
\endlastfoot
\lain{u8}          & \texttt{uint8\_t}   & \\
\lain{u16}         & \texttt{uint16\_t}  & \\
\lain{u32}         & \texttt{uint32\_t}  & \\
\lain{u64}         & \texttt{uint64\_t}  & \\
\lain{usize}       & \texttt{size\_t}    & pointer-sized \\
\lain{i8}          & \texttt{int8\_t}    & \\
\lain{i16}         & \texttt{int16\_t}   & \\
\lain{i32}         & \texttt{int32\_t}   & \\
\lain{i64}         & \texttt{int64\_t}   & \\
\lain{isize}       & \texttt{ptrdiff\_t} & pointer-sized \\
\lain{int}         & \texttt{int}        & at least 32-bit \\
\lain{f32}         & \texttt{float}      & IEEE 754 \\
\lain{f64}         & \texttt{double}     & IEEE 754 \\
\lain{bool}        & \texttt{\_Bool}     & C99 bool (1 byte); \idbehav{0 or 1} \\
shared struct \lain{T} & \texttt{const T*} & \\
\lain{var} struct \lain{T} & \texttt{T*}  & \\
\lain{mov} struct \lain{T} & \texttt{T}   & by value \\
\lain{T[]}         & \texttt{struct \{ size\_t len; T* data; \}} & fat slice \\
\lain{*T}          & \texttt{const T*}  & unsafe only \\
\lain{var *T}      & \texttt{T*}        & unsafe only \\
\end{longtable}

\section{ABI for function parameters}
\label{sec:interop.abi}
\index{ABI}

\pnum
This table is normative and applies to all Lain function calls, including
those to \lain{extern} functions.

\begin{longtable}{@{}llll@{}}
\toprule
\textbf{Lain param} & \textbf{Type} & \textbf{C param} & \textbf{Direction} \\
\midrule
\endfirsthead
\toprule
\textbf{Lain param} & \textbf{Type} & \textbf{C param} & \textbf{Direction} \\
\midrule
\endhead
\bottomrule
\endlastfoot
\lain{name T}     & struct & \texttt{const T*}  & in \\
\lain{var name T} & struct & \texttt{T*}        & in/out \\
\lain{mov name T} & struct & \texttt{T}         & in (by value) \\
any mode          & primitive & \texttt{T}      & in \\
\end{longtable}

\section{Pointer operations}
\label{sec:interop.pointers}
\index{pointer operations}

\pnum
Raw pointer operations (\lain{*ptr} dereference, \lain{\&x} address-of)
require an \lain{unsafe} block (see \S\,\ref{sec:18-unsafe}).

\begin{constraint}
  Pointer-to-pointer or integer-to-pointer casts shall occur only inside
  an \lain{unsafe} block~(\diagcode{E012}).
\end{constraint}

\section{String interop}
\label{sec:interop.strings}
\index{string interop}

\pnum
Lain strings (\lain{u8[]}) are byte slices, not null-terminated C strings.
To pass a Lain string to a C function expecting \texttt{char*}, the
\lain{.data} pointer must be extracted explicitly:

\begin{laincode}
c_include "<stdio.h>"
extern func puts(s *u8) int effects io

var s = "hello"
puts(s.data)   // OK: .data is *u8 = const char*
\end{laincode}

\pnum
String literals are null-terminated (\S\,\ref{sec:lex.literals.string}),
so \lain{literal.data} is a valid \texttt{const char*} for C functions
that expect a C string.

\section{Opaque C types}
\label{sec:interop.opaque}
\index{opaque type}

\pnum
An opaque C type is declared with \lain{extern type Name}.  The type's
layout is unknown to Lain; it may only be used behind a pointer.

\begin{laincode}
extern type FILE
extern func fopen(path *u8, mode *u8) mov *FILE effects io, alloc
\end{laincode}

\begin{constraint}
  An opaque type shall not be used as a value type (only as a pointer
  target).  \lain{var f FILE} is \illformed{}.
\end{constraint}
